\section{The finite obstruction on a cross-polytope}
\label{sec:finite-degree}

Write $d=n+1$. Let $K$ be the simplicial complex with vertices
$+1,-1,\ldots,+d,-d$ whose faces are the subsets containing no opposite
pair. Geometrically the signed index $\sigma i$ stands for $\sigma e_i$.
Thus $|K|=\{x\in\mathbb R^d:\sum_i|x_i|=1\}$, and every simplex used
below has rational vertices. Write $K_r$ for the $r$-fold barycentric
subdivision, with $K_0=K$. A vertex $x$ of $K_r$ has a nonempty signed
support, the signed indices of its nonzero coordinates. Put
\[
    c(x)=\operatorname{sgn}(x_i)i,\qquad
    i=\min\{j: |x_j|=\max_k |x_k|\}.
\]
In particular $c(x)$ belongs to the support of $x$.

\begin{theorem}[Finite obstruction]\label{thm:obstruction}
Suppose $d$ is odd, $r\geq0$, and $T\geq1$. At every vertex $x$ of
$K_r$ and every integer $j$ from $0$ through $T$, prescribe a label
$\ell_j(x)\in\{\pm1,\ldots,\pm d\}$, with
$\ell_0(x)=c(x)$ and $\ell_T(x)=-c(x)$. There is a simplex $s$ of
$K_r$ and a $j<T$ such that among the labels
\[
       \{\ell_j(x),\ell_{j+1}(x):x\text{ a vertex of }s\}
\]
one finds opposite signed indices. A bounded scan of the finite lists
returns $j$, $s$, two vertices and their layers witnessing this assertion.
\end{theorem}

\begin{proof}
We compare one oriented chain at the two ends of the proposed sequence.

\noindent\textit{Oriented chains.}
Use integer formal sums of oriented simplices. A tuple
$[v_0,\ldots,v_p]$ changes sign under a transposition and denotes zero
if two entries coincide. Its boundary for $p>0$ is the alternating sum
obtained by deleting one entry, and a vertex has zero boundary.
A vertex map $f$ taking each simplex into a simplex gives a chain map
$f_\#[v_0,\ldots,v_p]=[f(v_0),\ldots,f(v_p)]$. Indeed the
boundary of the tuple is obtained by deleting the same entries before
or after applying $f$. If the image has repeated entries, choose two
equal entries in positions $a<b$. All deletion terms except possibly
$a,b$ still have a repeat. The two exceptional lists differ by moving
one entry across $b-a-1$ others. Their orientation signs, combined
with their deletion signs, cancel.

For an oriented simplex $s=[v_0,\ldots,v_p]$, its oriented subdivision
$S(s)$ is the following chain. For each permutation $\pi$ of
$0,\ldots,p$ take, with coefficient $\operatorname{sgn}(\pi)$, the
flag of successive barycentres of
$\{v_{\pi(0)}\},\{v_{\pi(0)},v_{\pi(1)}\},\ldots,
\{v_{\pi(0)},\ldots,v_{\pi(p)}\}$, listed in that order. This formula
is alternating in the vertices of $s$. In its boundary, deletion of
any barycentre other than the last cancels with the flag obtained by
interchanging the two consecutive vertices added across the missing
stage. If the last vertex added was $v_i$, deletion of the last
barycentre has coefficient $(-1)^i$ times the permutation sign on
the other vertices. Hence $\partial S(s)=S(\partial s)$ for $p>0$,
and the same identity holds at $p=0$. Iterate this explicit operator
to obtain $S^r$. It sends any cycle to a cycle.

\par\smallskip\noindent\textit{What the initial labels do to subdivision.}\enspace
Every fine face inside a coarse face $s$ has all its vertices supported
on signed indices belonging to $s$. Consequently the labels $c$ on
that fine face belong to $s$. In fact
\[
                          c_\# S^r(s)=s
\]
for every oriented coarse face $s$. For a vertex the identity is
immediate. For a $p$-face with $p>0$ the left side is a $p$-chain
supported on that same face, hence an integer multiple $a s$.
The boundary identity and the assertion for the facets of $s$ give
$a\partial s=\partial s$. Comparing the coefficient of just one
facet yields $a=1$. This proves the identity by induction on $p$.

\par\smallskip\noindent\textit{The cycle distinguishing the endpoints.}\enspace
For a sign vector $\sigma=(\sigma_1,\ldots,\sigma_d)$ let
$F_\sigma=[\sigma_1 e_1,\ldots,\sigma_d e_d]$ in increasing
order of the unsigned indices. Define the top chain
\[
                    Z=\sum_{\sigma\in\{\pm1\}^d}
                           (\sigma_1\cdots\sigma_d) F_\sigma.
\]
It is a cycle. For $d\ge2$ a facet with the $i$th index missing occurs
twice with the same deletion sign and opposite coefficient. For $d=1$
the boundary of each vertex is zero. The antipodal map $A(v)=-v$
sends $F_\sigma$ to $F_{-\sigma}$. Reindexing the finite sum gives
$A_\#Z=(-1)^d Z$. The coefficient on the all-positive face of $Z$ is $1$.

\par\smallskip\noindent\textit{What adjacent layers preserve.}\enspace
Suppose the scan found no opposite pair. For each pair of consecutive
layers the union of the two label sets of a fine simplex is therefore
a face of $K$. Let $f=\ell_j$ and $g=\ell_{j+1}$, two simplicial
vertex maps. Fix a total order on the fine vertices, and on each
increasing oriented basis simplex put
\[
 P[v_0,\ldots,v_p]=\sum_{i=0}^p(-1)^i
 [f(v_0),\ldots,f(v_i),g(v_i),\ldots,g(v_p)] .
\]
Extend by integer linearity. The common-face property makes every
tuple legitimate. The identity
\[
                \partial P+P\partial=g_\#-f_\#
\]
is a finite cancellation. In summand $i$, deleting $f(v_i)$ matches
deleting $g(v_{i-1})$ in summand $i-1$ with opposite coefficient
when $i>0$. The unmatched end terms are $g_\#$ and $-f_\#$.
For a deletion $f(v_k)$ with $k<i$, its coefficient is $(-1)^{i+k}$,
opposite to the coefficient $(-1)^{k+i-1}$ in $P\partial$.
For a deletion $g(v_k)$ with $k>i$, its coefficient is
$(-1)^{i+k+1}$, opposite to $(-1)^{k+i}$ in $P\partial$.
These identities also apply when images repeat, by the alternating-tuple
rule above.

The chain $S^r(Z)$ has degree $n=d-1$. Each tuple in $P(S^r(Z))$
has $d+1$ entries belonging to a single face with at most $d$ vertices,
so this chain vanishes. Also $\partial S^r(Z)=0$. Thus consecutive
layers have exactly the same top image chain. Repeating this equality
gives $Z=c_\#S^r(Z)=(A\circ c)_\#S^r(Z)=(-1)^d Z$.
For odd $d$, the coefficient at the all-positive face would give
$1=-1$. Therefore the finite scan cannot end in failure. Its success
branch already records the required witnesses.
\end{proof}

\begin{lemma}[Mesh in the $\ell^1$ metric]\label{lem:mesh}
For $n\geq1$ the diameter of every simplex after $r=ns$ barycentric
subdivisions is at most $2^{1-s}$. For $n=0$ every diameter is zero.
\end{lemma}
\begin{proof}
A coarse face has diameter at most $2$ in the $\ell^1$ norm, since
every signed coordinate vertex has norm $1$. In a fine
simplex its barycentric vertices come from a nested family of faces.
For nonempty $B\supsetneq D$,
\[
 b_B=\frac{|D|}{|B|}b_D+
       \frac{|B|-|D|}{|B|}b_{B\setminus D}.
\]
The distance of these barycentres $b_B,b_D$ is at most
$n/(n+1)$ times the diameter of the parent simplex. Indeed the
second term uses a convex combination of vertices of that simplex,
and the same bound applies to every pair in the nested family.
For $n\geq1$, Bernoulli's inequality gives
$(1+1/n)^n\geq2$. After $r=ns$ subdivisions the diameter is at
most $2^{1-s}$. For $n=0$ the diameter is already zero.
\end{proof}

\par\smallskip\noindent\textit{The induction used by the obstruction.}\enspace
The finite assertion also has a uniform arithmetic reading with $n,r,T$
as numbers. An abstract simplex can be represented by an increasing list
of vertex identifiers, an oriented tuple by its sorted list and permutation
sign, and a chain by an array of signed integer coefficients indexed by
the enumerated faces of its degree. Store the coordinate tuple alongside
each vertex identifier. To form one subdivision, enumerate nonempty faces
of the preceding complex, assign each its average, then enumerate nested
lists of those faces; their lengths are at most $d$. For instance, from
$h$ face identifiers one may try all ordered lists of length up to $d$
over $0,\ldots,h-1$ and keep strictly nested lists. This step is a
bounded loop on explicitly given numbers, and repetition for $r$ steps
is ordinary primitive recursion. The signs in $S$, $P$ and $Z$ and all
chain sums similarly have primitive-recursive implementations. One can
use signed numerator/positive denominator arithmetic throughout.

In the carrier induction proving $c_\# S^r(s)=s$, the apparent assertion
that every coarse $p$-face satisfies the identity means the Boolean
conjunction of coefficient equalities over the already enumerated faces.
The induction on $p$ is therefore quantifier-free induction with $n,r$
as parameters. The same compilation applies to the boundary and mesh
identities at subdivision stages and to equality of the image chains at
time layers. In the sign cancellations, finite contributions of opposite
signs are paired by the explicit transposition or deletion rule just given.
Summing over the resulting pairs is bounded arithmetic, justified for
example by induction on the processed prefix of a finite list. Finally,
run the witness scan with a success flag. Its failure flag implies that
every tested pair passed the no-opposites check, whence the displayed
impossible integer equality. A case split on this one Boolean proves
success and returns the stored indices intuitionistically.
